For most deserts the only free elevation models a small wheeled rover can plan on have a posting of about 30\,m, while the main hazard on active dunes, the slip face, is often only metres long. We built a benchmark on six North American dune fields with one metre airborne lidar as truth. Ten families of hazard layers, from slope and interpolation to boosted trees, a convolutional network, learned super resolution and imagery fusion, were scored against a lidar face map and used by a slip aware planner at operating points chosen on the other five sites. At 10\,m the best layer raises mission success from {{VAL:nv_post_E1_I10}} to {{VAL:ho_post_E1_I10_BICUBICcal}}; at 30\,m no computable layer adds more than {{VAL:ho_post_E1_I30_DEMplusS2_dabs}} points to {{VAL:nv_post_E1_I30}}, even for faces longer than 60\,m, while a face map oracle reaches {{VAL:ho_post_E1_I30_FACEMAP}}, and fewer routes still are climbable on the lidar. In post hoc closed loop runs with line of sight sensing, 30\,m priors raise success from {{VAL:lo_I30_E1_none_p10}} to at most {{VAL:lo_I30_E1_BICUBICcal_p10}}, and not clearly when endpoints are chosen from the coarse map. Code, results and the prediction log are public.
